% Options for packages loaded elsewhere
\PassOptionsToPackage{unicode}{hyperref}
\PassOptionsToPackage{hyphens}{url}
\PassOptionsToPackage{dvipsnames,svgnames,x11names}{xcolor}
\documentclass[
  11pt,
  a4paper,
]{article}
\usepackage{xcolor}
\usepackage[margin=1in]{geometry}
\usepackage{amsmath,amssymb}
\setcounter{secnumdepth}{-\maxdimen} % remove section numbering
\usepackage{iftex}
\ifPDFTeX
  \usepackage[T1]{fontenc}
  \usepackage[utf8]{inputenc}
  \usepackage{textcomp} % provide euro and other symbols
\else % if luatex or xetex
  \usepackage{unicode-math} % this also loads fontspec
  \defaultfontfeatures{Scale=MatchLowercase}
  \defaultfontfeatures[\rmfamily]{Ligatures=TeX,Scale=1}
\fi
\usepackage{lmodern}
\ifPDFTeX\else
  % xetex/luatex font selection
\fi
% Use upquote if available, for straight quotes in verbatim environments
\IfFileExists{upquote.sty}{\usepackage{upquote}}{}
\IfFileExists{microtype.sty}{% use microtype if available
  \usepackage[]{microtype}
  \UseMicrotypeSet[protrusion]{basicmath} % disable protrusion for tt fonts
}{}
\usepackage{setspace}
\emergencystretch=3em
\sloppy
\newcommand{\codebreak}[1]{\begingroup\ttfamily\def\_{\textunderscore\allowbreak}#1\endgroup}
\makeatletter
\@ifundefined{KOMAClassName}{% if non-KOMA class
  \IfFileExists{parskip.sty}{%
    \usepackage{parskip}
  }{% else
    \setlength{\parindent}{0pt}
    \setlength{\parskip}{6pt plus 2pt minus 1pt}}
}{% if KOMA class
  \KOMAoptions{parskip=half}}
\makeatother
\usepackage{color}
\usepackage{fancyvrb}
\newcommand{\VerbBar}{|}
\newcommand{\VERB}{\Verb[commandchars=\\\{\}]}
\DefineVerbatimEnvironment{Highlighting}{Verbatim}{commandchars=\\\{\}}
% Add ',fontsize=\small' for more characters per line
\newenvironment{Shaded}{}{}
\newcommand{\AlertTok}[1]{\textcolor[rgb]{1.00,0.00,0.00}{\textbf{#1}}}
\newcommand{\AnnotationTok}[1]{\textcolor[rgb]{0.38,0.63,0.69}{\textbf{\textit{#1}}}}
\newcommand{\AttributeTok}[1]{\textcolor[rgb]{0.49,0.56,0.16}{#1}}
\newcommand{\BaseNTok}[1]{\textcolor[rgb]{0.25,0.63,0.44}{#1}}
\newcommand{\BuiltInTok}[1]{\textcolor[rgb]{0.00,0.50,0.00}{#1}}
\newcommand{\CharTok}[1]{\textcolor[rgb]{0.25,0.44,0.63}{#1}}
\newcommand{\CommentTok}[1]{\textcolor[rgb]{0.38,0.63,0.69}{\textit{#1}}}
\newcommand{\CommentVarTok}[1]{\textcolor[rgb]{0.38,0.63,0.69}{\textbf{\textit{#1}}}}
\newcommand{\ConstantTok}[1]{\textcolor[rgb]{0.53,0.00,0.00}{#1}}
\newcommand{\ControlFlowTok}[1]{\textcolor[rgb]{0.00,0.44,0.13}{\textbf{#1}}}
\newcommand{\DataTypeTok}[1]{\textcolor[rgb]{0.56,0.13,0.00}{#1}}
\newcommand{\DecValTok}[1]{\textcolor[rgb]{0.25,0.63,0.44}{#1}}
\newcommand{\DocumentationTok}[1]{\textcolor[rgb]{0.73,0.13,0.13}{\textit{#1}}}
\newcommand{\ErrorTok}[1]{\textcolor[rgb]{1.00,0.00,0.00}{\textbf{#1}}}
\newcommand{\ExtensionTok}[1]{#1}
\newcommand{\FloatTok}[1]{\textcolor[rgb]{0.25,0.63,0.44}{#1}}
\newcommand{\FunctionTok}[1]{\textcolor[rgb]{0.02,0.16,0.49}{#1}}
\newcommand{\ImportTok}[1]{\textcolor[rgb]{0.00,0.50,0.00}{\textbf{#1}}}
\newcommand{\InformationTok}[1]{\textcolor[rgb]{0.38,0.63,0.69}{\textbf{\textit{#1}}}}
\newcommand{\KeywordTok}[1]{\textcolor[rgb]{0.00,0.44,0.13}{\textbf{#1}}}
\newcommand{\NormalTok}[1]{#1}
\newcommand{\OperatorTok}[1]{\textcolor[rgb]{0.40,0.40,0.40}{#1}}
\newcommand{\OtherTok}[1]{\textcolor[rgb]{0.00,0.44,0.13}{#1}}
\newcommand{\PreprocessorTok}[1]{\textcolor[rgb]{0.74,0.48,0.00}{#1}}
\newcommand{\RegionMarkerTok}[1]{#1}
\newcommand{\SpecialCharTok}[1]{\textcolor[rgb]{0.25,0.44,0.63}{#1}}
\newcommand{\SpecialStringTok}[1]{\textcolor[rgb]{0.73,0.40,0.53}{#1}}
\newcommand{\StringTok}[1]{\textcolor[rgb]{0.25,0.44,0.63}{#1}}
\newcommand{\VariableTok}[1]{\textcolor[rgb]{0.10,0.09,0.49}{#1}}
\newcommand{\VerbatimStringTok}[1]{\textcolor[rgb]{0.25,0.44,0.63}{#1}}
\newcommand{\WarningTok}[1]{\textcolor[rgb]{0.38,0.63,0.69}{\textbf{\textit{#1}}}}
\usepackage{longtable,booktabs,array}
\newcounter{none} % for unnumbered tables
\usepackage{calc} % for calculating minipage widths
% Correct order of tables after \paragraph or \subparagraph
\usepackage{etoolbox}
\makeatletter
\patchcmd\longtable{\par}{\if@noskipsec\mbox{}\fi\par}{}{}
\makeatother
% Allow footnotes in longtable head/foot
\IfFileExists{footnotehyper.sty}{\usepackage{footnotehyper}}{\usepackage{footnote}}
\makesavenoteenv{longtable}
\setlength{\emergencystretch}{3em} % prevent overfull lines
\providecommand{\tightlist}{%
  \setlength{\itemsep}{0pt}\setlength{\parskip}{0pt}}
\usepackage{bookmark}
\IfFileExists{xurl.sty}{\usepackage{xurl}}{} % add URL line breaks if available
\urlstyle{same}
\hypersetup{
  pdftitle={Compressed GeoFM representations improve held-out farmland layout optimization under evidence gates},
  colorlinks=true,
  linkcolor={Maroon},
  filecolor={Maroon},
  citecolor={Blue},
  urlcolor={Blue},
  pdfcreator={LaTeX via pandoc}}

\title{Compressed GeoFM representations improve held-out farmland layout
optimization under evidence gates}
\author{}
\date{}

\begin{document}
\maketitle

\setstretch{1.15}
\section{Evidence-Gated GeoFM Representation for Farmland Layout
Optimization}\label{evidence-gated-geofm-representation-for-farmland-layout-optimization}

\subsection{Title}\label{title}

Compressed GeoFM representations improve held-out farmland layout
optimization under evidence gates

\subsection{Short Title}\label{short-title}

Compressed GeoFM planning

\subsection{Keywords}\label{keywords}

Geospatial foundation model; AlphaEarth embeddings; farmland layout
optimization; reinforcement learning; suitability proxy; representation
control; evidence gate; land-use planning

\subsection{Highlights}\label{highlights}

\begin{itemize}
\tightlist
\item
  A reproducible workflow links AlphaEarth embeddings to real Bishan
  DLTB planning blocks.
\item
  Tiled and padded policy interfaces make large real planning instances
  testable.
\item
  Multi-tile B0/B1 evidence rejects raw 64-dimensional GeoFM direct
  injection.
\item
  Phase 48 shows that compressed GeoFM routes D4P8 and D4P16 outperform
  B0, raw B1, random D2, and shuffled D3 on mean held-out reward.
\item
  Phase 52 expands this compressed-route evidence to five held-out tiles
  and three seeds.
\item
  Phase 53 supports the expanded cluster mean with exact sign-flip,
  bootstrap, and leave-one influence checks.
\item
  Phase 54 verifies that the formal Phase 52/53 values come from one
  authoritative artifact chain.
\item
  Suitability-reward work remains blocked until an independent
  non-leakage label gate passes.
\item
  A local Phase 42 label-source audit finds no usable external
  suitability label for the real Bishan run.
\item
  The evidence supports a bounded positive compressed-representation
  conclusion, not a raw-GeoFM or suitability-reward claim.
\end{itemize}

\subsection{Abstract}\label{abstract}

Farmland spatial layout optimization requires decisions over many
heterogeneous planning units, but direct variables for soil quality,
irrigation, productivity, and long-term suitability are often
unavailable at operational planning scale. Geospatial foundation-model
(GeoFM) embeddings offer a possible source of latent land-surface
information, but their usefulness for reinforcement-learning planning
depends on representation design, reward validation, and leakage-aware
labels. This study builds a reproducible evidence-gated workflow that
aggregates AlphaEarth embeddings and explicit planning features to real
Bishan DLTB land-use blocks, constructs tiled and padded planning
interfaces, and evaluates GeoFM-enhanced representations under a
deterministic base planning reward. Raw 64-dimensional GeoFM direct
injection is not supported: at 4096 training steps, B1 has a
learned-policy B1-B0 mean reward delta of \texttt{-0.1318712688}, with
only \texttt{3\ /\ 9} held-out tile-seed pairs favoring B1. Phase 48
changes the broader conclusion by re-evaluating the PCA-compressed GeoFM
variants as candidate state routes. D4P8 and D4P16 exceed B0, raw B1,
random D2, and shuffled D3 on mean learned-policy reward; the pooled
compressed-control delta is \texttt{0.4673011499}, with
\texttt{48\ /\ 72} positive comparisons. Phase 49 further reports
\codebreak{compressed\_route\_statistically\_robust}, with one-sided
sign-test p \texttt{0.0031549137} and bootstrap CI95
\texttt{{[}0.2827829983,\ 0.6639974489{]}}. Phase 50 aggregates the same
evidence to tile-seed clusters and reports directional support
(\texttt{7\ /\ 9} positive clusters, p \texttt{0.08984375}), so the
sign-only cluster test is underpowered. Phase 51 then uses an exact
signed-rank test over the cluster mean deltas and supports the
magnitude-sensitive cluster effect (positive rank sum
\texttt{40\ /\ 45}, p \texttt{0.01953125}). Phase 52 expands the
six-variant replication to five held-out tiles and three seeds; it again
reports \codebreak{compressed\_geofm\_route\_supported}, all eight
compressed-versus-control mean deltas positive, pooled delta
\texttt{0.2921767818} with \texttt{74\ /\ 120} positive comparisons,
row-level sign-test p \texttt{0.0066881634}, cluster sign-only p
\texttt{0.1508789062}, and cluster signed-rank p \texttt{0.0206298828}.
Phase 53 adds direct cluster-mean support: exact one-sided sign-flip
mean p \texttt{0.0196838379}, bootstrap CI95
\texttt{{[}0.0570820445,\ 0.5823557658{]}}, and positive leave-one
cluster, tile, and seed means. Phase 54 reports
\codebreak{artifact\_lineage\_consistent}, confirming that the formal Phase
52/53 values are reproducible from one authoritative
delta-to-cluster-to-statistics artifact chain. D4P16 reaches mean reward
\texttt{0.9918299718} in the original three-tile Phase 48 audit and
\texttt{0.5819662325} in the expanded Phase 52 audit. Normalized-B1 and
budget checks do not displace the compressed route: Phase 30 improves
raw B1 but remains below D4P8/D4P16, and Phase 33 reports
\codebreak{budget\_not\_explanatory} for the normalized branch.
Suitability-reward evidence remains blocked: weak labels are
DLTB/slope-derived leakage risks, the scalar \texttt{suitability\_proxy}
is not supported as a reward term, and Phase 40/41 report missing
independent label inputs. These results support a bounded positive
conclusion for compressed GeoFM state representations under the current
Bishan base-reward held-out protocol, while raw direct injection, B2/B3
suitability reward, and cross-region transfer claims remain unsupported.

\subsection{1. Introduction}\label{introduction}

Farmland layout optimization is not only a geometric consolidation
problem. In practical land-use planning, decisions about whether to
retain, exchange, or consolidate farmland depend on explicit spatial
constraints as well as latent environmental conditions that are
difficult to observe comprehensively at the planning-unit scale. Slope,
contiguity, block area, adjacency, and fragmentation can often be
derived from GIS data, but variables related to irrigation, soil
quality, moisture, and productivity may be incomplete, unavailable, or
restricted.

Geospatial foundation-model embeddings provide a tempting route around
this data gap. Multi-sensor satellite embedding products such as
AlphaEarth can encode broad land-surface context in compact feature
vectors, and such vectors may contain information relevant to farmland
suitability. However, a remote- sensing embedding is not a direct
measurement of agronomic suitability. If a planning model improves after
adding GeoFM embeddings, the improvement may come from meaningful
environmental signal, from representation scale, from redundant
compression, from easier optimization geometry, or from leakage through
labels or features.

Paper11 was designed to test this distinction rather than assume it. The
workflow links AlphaEarth embeddings to real Bishan DLTB land-use
polygons, constructs explicit-only and GeoFM-enhanced planning
representations, and evaluates learned policies under a deterministic
base planning reward. It also implements random, shuffled, compressed,
normalized, suitability-proxy, and independent-label gates so that each
stronger manuscript claim must pass an explicit evidence threshold.

The current evidence refines rather than rejects the original positive
framing. Raw GeoFM direct injection is not a supported improvement, and
suitability reward must not be used without independent labels. However,
Phase 48 shows that compressed GeoFM state routes are supported under
the current base-reward held-out protocol. This paper therefore frames
Paper11 as a bounded positive compressed-representation result with
explicit evidence gates around raw B1, suitability reward, and transfer
claims.

\subsection{2. Data and Planning Units}\label{data-and-planning-units}

The workflow uses real Bishan DLTB land-use polygons as planning units.
The Phase 11 adapter exports \texttt{64,984} DLTB polygons into Phase
2-compatible block feature tables. The feature pipeline aggregates
AlphaEarth annual embeddings to block-level representations and joins
explicit planning attributes such as slope- and land-use-derived
variables where available.

The repository includes lightweight sample AlphaEarth arrays for smoke
tests, while the full external Bishan DLTB-with-slope GeoPackage is
documented as a local external input and is not redistributed in
ordinary Git. Generated real outputs under
\texttt{experiments/**/outputs/} are ignored artifacts. The repository
therefore separates lightweight reviewer reproducibility from full local
data reconstruction.

\subsection{3. Method}\label{method}

\subsubsection{3.1 Representation
Families}\label{representation-families}

The core representation families are:

\begin{itemize}
\tightlist
\item
  \texttt{B0}: explicit planning features with deterministic
  \texttt{base\_planning\_reward};
\item
  \texttt{B1}: explicit planning features plus raw 64-dimensional
  AlphaEarth or GeoFM embeddings;
\item
  \texttt{D2}: explicit planning features plus random 64-dimensional
  controls;
\item
  \texttt{D3}: explicit planning features plus shuffled AlphaEarth
  embeddings;
\item
  \texttt{D4P8} and \texttt{D4P16}: explicit planning features plus
  PCA-compressed AlphaEarth embeddings;
\item
  \texttt{N1Z} and \texttt{N1ZR}: normalized B1 variants used to test
  representation-scale effects;
\item
  \texttt{B2} and \texttt{B3}: suitability-reward variants reserved for
  later use and not enabled in the current evidence.
\end{itemize}

This design separates three questions. First, can the planning
environment and policy interface operate on real planning units? Second,
does GeoFM information improve learned-policy decisions under the same
deterministic base reward, and which state representation makes that
information usable? Third, is there validated suitability evidence
strong enough to justify adding a suitability reward? The current study
answers the first question positively, the second positively only for
compressed GeoFM state routes, and the third negatively under present
inputs.

\subsubsection{3.2 Planning Reward and Policy
Interface}\label{planning-reward-and-policy-interface}

The current learned-policy experiments use a deterministic
\texttt{base\_planning\_reward}. This reward encodes explicit planning
logic such as slope, contiguity, baimu-fang-style consolidation, action
validity, and related planning constraints. It deliberately excludes
suitability reward because the suitability proxy has not passed the
evidence gates required for reward integration.

Early flat observation/action interfaces were tile-size-specific and
could not support learned-policy evaluation across variable-size
held-out tiles. Later phases introduced tiled contracts and padded
variable-size policy evaluation, allowing bounded same-reward
comparisons across held-out Bishan tiles.

\subsubsection{3.3 Evidence Gates}\label{evidence-gates}

The evidence-gated design prevents unsupported claims from entering the
manuscript. Phase 26 converts padded held-out policy outputs into
manuscript-facing B0/B1 summaries. Phase 28 adds representation
controls. Phase 33 tests whether a modestly higher budget rescues
normalized-B1 behavior. Phase 48 re-evaluates D4P8 and D4P16 as
compressed GeoFM candidate routes under the same held-out base-reward
protocol. Phase 49-54 then test statistical robustness, tile-seed
clustering, signed-rank cluster magnitude support, an expanded
five-tile, three-seed replication, cluster-mean influence support, and
artifact-lineage consistency for the same compressed route. Phase 36
evaluates suitability-proxy signal against available weak labels. Phase
38 rebuilds suitability proxies under leakage-aware controls. Phase 39
audits available label-like columns. Phase 40 adds the hard
independent-label go/no-go gate before any Phase 38 rerun or B2/B3
reward smoke. Phase 41 tests a separate suitability-prior route in which
GeoFM can enter reward design only after independent-label calibration.
Phase 42 audits local candidate label sources to distinguish true
external labels from diagnostic DLTB/slope weak labels and unrelated
labels from other projects.

\subsection{4. Experiments}\label{experiments}

\subsubsection{4.1 Held-Out B0/B1 Learned-Policy
Analysis}\label{held-out-b0b1-learned-policy-analysis}

The main B0/B1 held-out analysis trains on \texttt{tile\_r003\_c003} and
evaluates on held-out Bishan tiles \texttt{tile\_r002\_c003},
\texttt{tile\_r005\_c004}, and \texttt{tile\_r005\_c003} across seeds
\texttt{0}, \texttt{1}, and \texttt{2}. At 4096 training steps and
evaluation horizon \texttt{8}, Phase 26 reports a learned-policy B1-B0
mean reward delta of \texttt{-0.1318712688}. Only \texttt{3\ /\ 9}
tile-seed pairs favor B1.

This result does not support a stable claim that raw GeoFM-enhanced B1
outperforms explicit-feature B0 under the deterministic base planning
reward.

\subsubsection{4.2 Raw and Compressed Representation-Control
Results}\label{raw-and-compressed-representation-control-results}

Phase 28 first compares B1 against B0, D2, D3, D4P8, and D4P16 under the
same padded held-out protocol at 1024 and 4096 training steps. Both runs
report \codebreak{compression\_matches\_raw} rather than
\texttt{representation\_signal\_supported} for raw B1. At 4096 steps,
the B1-comparator mean deltas are:

{\def\LTcaptype{none} % do not increment counter
\begin{longtable}[]{@{}
  >{\raggedright\arraybackslash}p{(\linewidth - 4\tabcolsep) * \real{0.2727}}
  >{\raggedleft\arraybackslash}p{(\linewidth - 4\tabcolsep) * \real{0.3636}}
  >{\raggedleft\arraybackslash}p{(\linewidth - 4\tabcolsep) * \real{0.3636}}@{}}
\toprule\noalign{}
\begin{minipage}[b]{\linewidth}\raggedright
Comparison
\end{minipage} & \begin{minipage}[b]{\linewidth}\raggedleft
B1 minus comparator mean delta
\end{minipage} & \begin{minipage}[b]{\linewidth}\raggedleft
Positive tile-seed count
\end{minipage} \\
\midrule\noalign{}
\endhead
\bottomrule\noalign{}
\endlastfoot
B1 - B0 & \texttt{-0.1318712688} & \texttt{3\ /\ 9} \\
B1 - D2 & \texttt{0.0744591656} & \texttt{4\ /\ 9} \\
B1 - D3 & \texttt{-0.1094750135} & \texttt{3\ /\ 9} \\
B1 - D4P8 & \texttt{-0.3768518347} & \texttt{2\ /\ 9} \\
B1 - D4P16 & \texttt{-0.6411940236} & \texttt{2\ /\ 9} \\
\end{longtable}
}

This result rejects raw 64-dimensional GeoFM direct injection as the
effective state route. It does not reject GeoFM information itself,
because the compressed GeoFM variants were the strongest variants in the
same held-out summary set.

Phase 48 therefore re-evaluates D4P8 and D4P16 as compressed GeoFM
candidate routes. The audit is read-only over the frozen 4096-step Phase
28 summary CSV and compares each compressed candidate against B0, raw
B1, random D2, and shuffled D3 on the same tile-seed pairs.

{\def\LTcaptype{none} % do not increment counter
\begin{longtable}[]{@{}
  >{\raggedright\arraybackslash}p{(\linewidth - 4\tabcolsep) * \real{0.2727}}
  >{\raggedleft\arraybackslash}p{(\linewidth - 4\tabcolsep) * \real{0.3636}}
  >{\raggedleft\arraybackslash}p{(\linewidth - 4\tabcolsep) * \real{0.3636}}@{}}
\toprule\noalign{}
\begin{minipage}[b]{\linewidth}\raggedright
Comparison
\end{minipage} & \begin{minipage}[b]{\linewidth}\raggedleft
Compressed minus comparator mean delta
\end{minipage} & \begin{minipage}[b]{\linewidth}\raggedleft
Positive tile-seed count
\end{minipage} \\
\midrule\noalign{}
\endhead
\bottomrule\noalign{}
\endlastfoot
D4P8 - B0 & \texttt{0.2449805659} & \texttt{4\ /\ 9} \\
D4P8 - B1 & \texttt{0.3768518347} & \texttt{7\ /\ 9} \\
D4P8 - D2 & \texttt{0.4513110003} & \texttt{7\ /\ 9} \\
D4P8 - D3 & \texttt{0.2673768211} & \texttt{6\ /\ 9} \\
D4P16 - B0 & \texttt{0.5093227548} & \texttt{5\ /\ 9} \\
D4P16 - B1 & \texttt{0.6411940236} & \texttt{7\ /\ 9} \\
D4P16 - D2 & \texttt{0.7156531892} & \texttt{5\ /\ 9} \\
D4P16 - D3 & \texttt{0.5317190100} & \texttt{7\ /\ 9} \\
\end{longtable}
}

Phase 48 reports \codebreak{compressed\_geofm\_route\_supported}. The
pooled compressed-control delta is \texttt{0.4673011499}, with
\texttt{48\ /\ 72} positive comparisons. The supported representation
claim is therefore compressed and bounded: D4P8 and D4P16 improve mean
learned-policy reward under the current Bishan base-reward held-out
protocol, while raw B1 remains unsupported. Phase 49 strengthens this
result with a pooled one-sided sign-test p of \texttt{0.0031549137},
bootstrap CI95 \texttt{{[}0.2827829983,\ 0.6639974489{]}}, and positive
leave-one-tile and leave-one-seed sensitivity checks. Phase 50 then
aggregates to \texttt{9} tile-seed clusters and reports directional
support (\texttt{7\ /\ 9} positive, p \texttt{0.08984375}), which
narrows sign-test wording. Phase 51 uses exact signed-rank evidence over
cluster magnitudes and supports the cluster-level effect with p
\texttt{0.01953125}.

Phase 52 expands this replication to five held-out tiles and three seeds
while keeping the same six variants, \texttt{4096} training steps, and
evaluation horizon \texttt{8}. The expanded run again reports
\codebreak{compressed\_geofm\_route\_supported}: all eight
compressed-versus-control mean deltas remain positive, and the pooled
compressed-control delta is \texttt{0.2921767818} with
\texttt{74\ /\ 120} positive row-level comparisons.

{\def\LTcaptype{none} % do not increment counter
\begin{longtable}[]{@{}
  >{\raggedright\arraybackslash}p{(\linewidth - 4\tabcolsep) * \real{0.2727}}
  >{\raggedleft\arraybackslash}p{(\linewidth - 4\tabcolsep) * \real{0.3636}}
  >{\raggedleft\arraybackslash}p{(\linewidth - 4\tabcolsep) * \real{0.3636}}@{}}
\toprule\noalign{}
\begin{minipage}[b]{\linewidth}\raggedright
Expanded Phase 52 comparison
\end{minipage} & \begin{minipage}[b]{\linewidth}\raggedleft
Compressed minus comparator mean delta
\end{minipage} & \begin{minipage}[b]{\linewidth}\raggedleft
Positive tile-seed count
\end{minipage} \\
\midrule\noalign{}
\endhead
\bottomrule\noalign{}
\endlastfoot
D4P8 - B0 & \texttt{0.2896842037} & \texttt{9\ /\ 15} \\
D4P8 - B1 & \texttt{0.2050431914} & \texttt{9\ /\ 15} \\
D4P8 - D2 & \texttt{0.2506866266} & \texttt{10\ /\ 15} \\
D4P8 - D3 & \texttt{0.1973780839} & \texttt{10\ /\ 15} \\
D4P16 - B0 & \texttt{0.4026417146} & \texttt{8\ /\ 15} \\
D4P16 - B1 & \texttt{0.3180007023} & \texttt{10\ /\ 15} \\
D4P16 - D2 & \texttt{0.3636441375} & \texttt{7\ /\ 15} \\
D4P16 - D3 & \texttt{0.3103355948} & \texttt{11\ /\ 15} \\
\end{longtable}
}

The expanded Phase 49-style robustness check remains
\codebreak{compressed\_route\_statistically\_robust}, with pooled row-level
sign-test p \texttt{0.0066881634}, bootstrap CI95
\texttt{{[}0.1623326461,\ 0.4323997354{]}}, and positive leave-one-tile
and leave-one-seed means. The expanded tile-seed sign-only test remains
directional (\texttt{10\ /\ 15} positive clusters, p
\texttt{0.1508789062}), but the exact signed-rank test over cluster
magnitudes remains significant (positive rank sum \texttt{96\ /\ 120}, p
\texttt{0.0206298828}).

Phase 53 audits whether the expanded cluster mean is driven only by one
or two large positive clusters. The status is
\codebreak{cluster\_mean\_support}: the cluster mean delta remains
\texttt{0.2921767818}, the exact one-sided sign-flip mean p is
\texttt{0.0196838379}, the bootstrap CI95 is
\texttt{{[}0.0570820445,\ 0.5823557658{]}}, and the minimum
leave-one-cluster, leave-one-tile, and leave-one-seed means are
\texttt{0.2060081575}, \texttt{0.0954244478}, and \texttt{0.2083797951},
respectively.

Phase 54 then verifies the artifact lineage for the formal Phase 52/53
evidence chain. It reports \codebreak{artifact\_lineage\_consistent}:
recomputing Phase 50 cluster means from the authoritative Phase 48 delta
table and recomputing Phase 51 and Phase 53 statistics from the
authoritative cluster CSV reproduces the values used in this manuscript.

\subsubsection{4.3 Normalization and Budget
Robustness}\label{normalization-and-budget-robustness}

Phase 30 tests whether normalizing B1 embeddings repairs raw-B1
underperformance. At 4096 steps, normalized variants improve raw B1 and
recover the mean B0 gap: \texttt{N1Z} achieves mean learned-policy
reward \texttt{0.6515323140}, compared with \texttt{0.4825072170} for B0
and \texttt{0.3506359482} for raw B1. However, normalized variants do
not consistently exceed compressed controls.

Phase 33 extends this question with bounded 5120-step matched pilots
across three held-out tiles and three seeds. The full aggregate reports
\codebreak{budget\_not\_explanatory}. At 5120 steps, tracked focal gaps
remain negative on mean, including
\texttt{N1Z\ -\ D4P16\ =\ -0.7597285327},
\texttt{N1Z\ -\ D4P8\ =\ -0.4953863438},
\texttt{N1ZR\ -\ D4P16\ =\ -0.6658915329}, and
\texttt{N1ZR\ -\ D4P8\ =\ -0.4015493440}.

This result prevents the manuscript from claiming that a modestly higher
training budget or simple normalization resolves the
representation-control problem.

\subsubsection{4.4 Suitability-Proxy, Prior, and Independent-Label
Gates}\label{suitability-proxy-prior-and-independent-label-gates}

Phase 36 evaluates whether GeoFM-derived feature families add weak-label
suitability signal beyond explicit planning features. The available
labels are \texttt{current\_farmland\_label},
\texttt{farmland\_or\_orchard\_label}, and
\texttt{low\_slope\_farmland\_label}, with \texttt{64,984} valid labels
each. All three labels are DLTB/slope-derived and flagged as
\texttt{explicit\_label\_leakage\_risk}.

Because explicit planning features encode the weak-label logic,
\texttt{explicit\_only} models reach ROC AUC, average precision, and
balanced accuracy of \texttt{1.0} for all three labels. This perfect
performance is a leakage warning, not suitability validation. GeoFM-only
checks are weak: for example, raw GeoFM-only features reach ROC AUC
\texttt{0.6490064144} for \texttt{low\_slope\_farmland\_label}, but the
scalar \texttt{suitability\_proxy} is near random for that label with
ROC AUC \texttt{0.4979564572}.

Phase 38 rebuilds proxy classifiers but remains
\codebreak{proxy\_rebuild\_diagnostic\_only} because evaluated labels are
leakage risks or GeoFM-derived proxies do not clear control thresholds.
Phase 39 reports \codebreak{independent\_label\_inputs\_missing} after
auditing the real Phase 2 table. Phase 40 then formalizes the decision
as a hard gate: without a registered independent non-leakage label,
Phase 38 cannot be rerun as a stronger proxy validation, and B2/B3
suitability-reward experiments should not proceed. The current Phase 40
no-registry run reads \texttt{64,984} feature rows, \texttt{0} registry
rows, and reports \codebreak{independent\_label\_inputs\_missing}.

Phase 41 tests a separate suitability route for GeoFM: after the
compressed state route has been evaluated under the base reward, any
suitability-prior export still requires independent-label calibration
and must pass explicit-baseline, shuffled-control, random-control,
fold-stability, and calibration checks. The current real no-registry run
reports \codebreak{phase41\_independent\_label\_inputs\_missing}, reads
\texttt{64,984} feature rows, finds \texttt{0} Phase-40-passed labels,
and produces no \codebreak{block\_geofm\_suitability\_prior.csv}.

Phase 42 audits local candidate label sources after Phase 41. DLTB/slope
labels can be registered only as diagnostic leakage-risk labels: Phase
40 reports \codebreak{independent\_label\_gate\_diagnostic\_only}, and
Phase 41 still reports
\codebreak{phase41\_independent\_label\_inputs\_missing} with the same
diagnostic registry. No local external soil, irrigation, yield,
productivity, high-standard-farmland, field-survey, or policy-outcome
label was found that can pass Phase 40.

\subsection{5. Discussion}\label{discussion}

The current Paper11 evidence supports a bounded positive representation
claim. The project began with the question of whether GeoFM embeddings
could improve farmland layout optimization by adding latent
environmental context. The answer is conditional. Raw 64-dimensional B1
direct injection does not stably outperform B0 across held-out Bishan
tiles and seeds. However, Phase 48 shows that compressed GeoFM state
routes D4P8 and D4P16 outperform B0, raw B1, random D2, and shuffled D3
on mean learned-policy reward under the same base-reward held-out
protocol. Phase 49 shows row-level statistical robustness within the
current Bishan protocol, while Phase 50 shows directional support after
conservative tile-seed cluster aggregation; Phase 51 adds exact
signed-rank cluster magnitude support. Phase 52 then expands the same
six-variant protocol to five held-out tiles and three seeds and
preserves the same conclusion: all compressed-versus-control mean deltas
remain positive, row-level evidence is statistically robust, and cluster
magnitude remains significant even though cluster signs alone remain
underpowered. Phase 53 adds that the expanded cluster mean itself is
supported by exact sign-flip, bootstrap, and leave-one influence checks,
so the positive compressed-route conclusion is not driven only by one
favorable cluster, tile, or seed. Phase 54 verifies that these formal
Phase 52/53 values are internally reproducible from one authoritative
artifact chain, preventing mixed-output lineage from supporting the
manuscript claim.

This distinction matters scientifically. A negative raw-B1 result could
have been misread as evidence that GeoFM information is irrelevant to
farmland layout optimization. The compressed-route result instead points
to representation geometry as the limiting factor. In the current
planning environment, GeoFM information appears useful when it enters
the policy state through compact PCA-compressed inputs, but not when raw
low-variance embedding dimensions are appended directly to explicit
planning features.

The result remains bounded. The Phase 48 audit is read-only over the
existing 4096-step Bishan held-out summary rows, and Phase 52 expands
the same protocol within Bishan rather than testing a new region. These
analyses do not prove that PCA is intrinsically optimal, that the effect
transfers across regions, or that the policy has learned agronomic
suitability semantics. The positive claim is about state representation
under a deterministic base planning reward, not about a validated
suitability reward.

The suitability branch is still more constrained than a simple missing
experiment. The current weak labels are derived from DLTB, slope, or
source metadata, so they cannot validate an independent suitability
reward. Phase 40 prevents the workflow from converting weak,
leakage-prone labels into B2/B3 reward claims. Phase 41 defines a
stricter future route in which GeoFM must be converted into a calibrated
low-dimensional prior before it can influence suitability reward. Phase
42 then tests whether the required labels are already available locally
and concludes that they are not: DLTB/slope labels remain
diagnostic-only, and candidate labels from Paper10 and Paper58 are not
Paper11 Bishan suitability labels.

The practical implication is that Paper11 should proceed with two
separate claim levels. The base-reward representation claim is now
supported for compressed GeoFM state routes in Bishan. The
suitability-reward claim remains blocked until external non-leakage
labels pass the independent-label gate and a leakage-aware GeoFM prior
clears control and calibration checks.

\subsection{6. Conclusion}\label{conclusion}

The current Paper11 repository establishes a reproducible workflow for
testing GeoFM-enhanced farmland layout optimization on real planning
units. The completed evidence supports a compressed GeoFM representation
route under the Bishan base-reward held-out protocol. Raw B1 direct
injection remains unsupported, but D4P8 and D4P16 improve over B0, raw
B1, random D2, and shuffled D3 on mean learned-policy reward, Phase 49
robustness checks keep the pooled effect positive, and Phase 50 keeps
the cluster-level direction positive with sign-only p
\texttt{0.08984375}; Phase 51 supports the cluster magnitude effect with
signed-rank p \texttt{0.01953125}. Phase 52 expands the same six-variant
test to five held-out tiles and three seeds, again supporting the
compressed route with pooled delta \texttt{0.2921767818},
\texttt{74\ /\ 120} positive row-level comparisons, row-level sign-test
p \texttt{0.0066881634}, and cluster signed-rank p
\texttt{0.0206298828}. Phase 53 further supports the expanded cluster
mean with exact sign-flip p \texttt{0.0196838379}, bootstrap CI95
\texttt{{[}0.0570820445,\ 0.5823557658{]}}, and positive leave-one
cluster, tile, and seed means. Phase 54 verifies the formal artifact
lineage as \codebreak{artifact\_lineage\_consistent}. The appropriate
current conclusion is therefore not that GeoFM fails, but that GeoFM
must be represented through a controlled compressed state route before
it improves the learned planning policy in this setting.
Suitability-reward work remains blocked until an independent non-leakage
label registry passes Phase 40 and a leakage-aware GeoFM prior passes
Phase 41.

\subsection{Claim-Evidence Map}\label{claim-evidence-map}

{\def\LTcaptype{none} % do not increment counter
\begin{longtable}[]{@{}
  >{\raggedright\arraybackslash}p{(\linewidth - 4\tabcolsep) * \real{0.3333}}
  >{\raggedright\arraybackslash}p{(\linewidth - 4\tabcolsep) * \real{0.3333}}
  >{\raggedright\arraybackslash}p{(\linewidth - 4\tabcolsep) * \real{0.3333}}@{}}
\toprule\noalign{}
\begin{minipage}[b]{\linewidth}\raggedright
Claim
\end{minipage} & \begin{minipage}[b]{\linewidth}\raggedright
Evidence
\end{minipage} & \begin{minipage}[b]{\linewidth}\raggedright
Status
\end{minipage} \\
\midrule\noalign{}
\endhead
\bottomrule\noalign{}
\endlastfoot
The repository implements a reproducible GeoFM-enhanced farmland
planning workflow. & Phase 1-25 pipeline, real Bishan DLTB adapter,
tiled contracts, padded held-out policy runner, smoke tests. &
Supported \\
Raw GeoFM B1 improves learned-policy planning decisions. & Phase 26
B1-B0 mean delta \texttt{-0.1318712688}, \texttt{3\ /\ 9} positive
tile-seed pairs. & Not supported \\
Raw B1 carries a stable representation advantage over controls. & Phase
28 reports \codebreak{compression\_matches\_raw}; D4P8 and D4P16 exceed B1
at 4096 steps. & Not supported \\
Compressed GeoFM state routes improve learned-policy planning decisions
under the base reward. & Phase 48 reports
\codebreak{compressed\_geofm\_route\_supported}; pooled compressed-control
delta \texttt{0.4673011499}, \texttt{48\ /\ 72} positive comparisons.
Phase 52 expands the six-variant run to five held-out tiles and three
seeds with pooled delta \texttt{0.2921767818} and \texttt{74\ /\ 120}
positive comparisons. & Supported \\
The compressed GeoFM state-route effect is row-level statistically
robust within the current Bishan protocol. & Phase 49 reports
\codebreak{compressed\_route\_statistically\_robust}; sign-test p
\texttt{0.0031549137}; bootstrap CI95
\texttt{{[}0.2827829983,\ 0.6639974489{]}}; leave-one tile/seed means
remain positive. Phase 52 repeats row-level robustness with sign-test p
\texttt{0.0066881634} and bootstrap CI95
\texttt{{[}0.1623326461,\ 0.4323997354{]}}. & Supported \\
The compressed GeoFM route is supported after tile-seed cluster
aggregation when magnitude and mean support are considered. & Phase 50
reports sign-only \codebreak{cluster\_directional\_support}; Phase 51
reports \codebreak{cluster\_magnitude\_support}, positive rank sum
\texttt{40\ /\ 45}, signed-rank p \texttt{0.01953125}. Phase 52 keeps
sign-only cluster support directional (\texttt{10\ /\ 15}, p
\texttt{0.1508789062}) but supports cluster magnitude with positive rank
sum \texttt{96\ /\ 120}, signed-rank p \texttt{0.0206298828}. Phase 53
reports \codebreak{cluster\_mean\_support}, exact sign-flip p
\texttt{0.0196838379}, bootstrap CI95
\texttt{{[}0.0570820445,\ 0.5823557658{]}}, and positive leave-one
cluster/tile/seed means. & Supported with magnitude-sensitive and
cluster-mean checks \\
The formal Phase 52/53 artifact chain is internally reproducible. &
Phase 54 reports \codebreak{artifact\_lineage\_consistent}: Phase 50
cluster means recomputed from the authoritative Phase 48 delta table,
and Phase 51/53 statistics recomputed from the authoritative cluster
CSV, match the formal manuscript values. & Supported \\
Normalization or more budget resolves the raw-B1 representation problem.
& Phase 30 partially improves B1; Phase 33 full bounded aggregate
reports \codebreak{budget\_not\_explanatory}. & Not supported \\
Suitability reward is ready for B2/B3. & Phase 36
\codebreak{proxy\_signal\_not\_supported}; Phase 38
\codebreak{proxy\_rebuild\_diagnostic\_only}; Phase 40
\codebreak{independent\_label\_inputs\_missing}; Phase 41
\codebreak{phase41\_independent\_label\_inputs\_missing}. & Not
supported \\
Local files already contain a usable independent Paper11 suitability
label. & Phase 42 finds only diagnostic DLTB/slope weak labels and
unrelated Paper10/Paper58 labels. & Not supported \\
The evidence supports a broad GeoFM-performance or transfer conclusion.
& Suitability, B2/B3, and transfer claims remain unsupported. & Not
supported \\
The evidence supports a bounded positive compressed-GeoFM representation
conclusion. & D4P8/D4P16 exceed B0, raw B1, D2, and D3 on mean held-out
base-reward policy reward; Phase 53 supports the expanded cluster mean
and influence checks. & Supported \\
\end{longtable}
}

\subsection{Data Availability}\label{data-availability}

The repository includes lightweight Bishan AlphaEarth sample arrays,
code, tests, reproduction guides, and file manifests for reviewer smoke
checks. Full real Bishan DLTB-with-slope inputs are external local data
and are not redistributed in ordinary Git. Large derived arrays, full
generated outputs, and trained artifacts should be deposited in an
external archive before any submission that relies on them.

\subsection{Code Availability}\label{code-availability}

All code for the reviewer-facing workflow is maintained in the Paper11
repository:

\begin{Shaded}
\begin{Highlighting}[]
\NormalTok{https://github.com/zhouning/paper11{-}geofm{-}farmland{-}suitability{-}rl}
\end{Highlighting}
\end{Shaded}

The final submitted version should cite a release tag or immutable
commit hash.

\subsection{Submission Metadata Still
Required}\label{submission-metadata-still-required}

\begin{itemize}
\tightlist
\item
  Final author list, affiliations, and corresponding-author details.
\item
  Funding statement.
\item
  Author contributions.
\item
  Final data access statement for the external DLTB-with-slope
  GeoPackage.
\item
  Reference list and citation formatting for the selected journal.
\item
  Final decision on target journal format for a bounded compressed-GeoFM
  representation article.
\end{itemize}

\end{document}
